\section{Models and Training}
\label{sec:models}

All models take a sequence of \V{data.frames} frames and output \V{data.signs} logits. The
\emph{GRU} runs a two-layer bidirectional GRU \cite{cho2014gru}; the \emph{CNN} is a stack of
residual depthwise-separable temporal convolutions \cite{howard2017mobilenets,bai2018tcn}; the
\emph{Transformer} is a four-layer encoder \cite{vaswani2017attention}. The \emph{hybrid}, close
to the Conformer \cite{gulati2020conformer}, has three stages of two depthwise-convolution blocks
and one transformer block, then average pooling over time and a linear head. Velocity, the v2
distance features and standardization are computed inside the network, so the exported ONNX file
takes raw normalized landmarks. The \emph{v1 recipe} trains for up to \V{hp.v1.epochs} epochs
with AdamW, cosine decay, label smoothing and geometric and temporal augmentation. The \emph{v2
recipe} trains up to \V{hp.v2.epochs} epochs at a lower learning rate, with an exponential moving
average of the weights \cite{polyak1992averaging}, frame dropout, time and landmark masking and
hand rotation. Each configuration has \V{hp.seeds} seeds; a hybrid run took about
\V{m.d.trainmin} minutes on a Kaggle \V{m.gpu}. The v2 grid is a cumulative ablation from the v1
CNN: v2 features, then the v2 recipe, then the hybrid.

% Generated by scripts/paper_tables.py -- do not edit.
\typeout{[table models]}
\begin{table*}[t]
\caption{Signer-independent models on the Google ISLR split}
\label{tab:models}
\centering
\small
\setlength{\tabcolsep}{5pt}
\begin{tabular}{lllrrrrrr}
\toprule
Model & Features & Recipe & Params (M) & Val top-1 & Test top-1 & Test top-3 & ONNX (MB) & Browser p95 (ms) \\
\midrule
GRU & v1 & v1 & 0.71 & 61.0 & 61.0$\pm$0.3 & 78.5$\pm$0.3 & 2.9 & 4.2 \\
Transformer & v1 & v1 & 1.28 & 62.1 & 62.5$\pm$0.2 & 78.0$\pm$0.3 & 5.2 & 9.1 \\
CNN & v1 & v1 & 0.32 & 62.2 & 62.5$\pm$0.2 & 79.5$\pm$0.5 & 1.3 & 3.3 \\
CNN & v2 & v1 & 0.40 & 61.0 & 62.7$\pm$0.3 & 80.5$\pm$0.1 & 1.6 & 4.8 \\
CNN & v2 & v2 & 0.40 & 62.4 & 63.7$\pm$0.0 & 81.0$\pm$0.2 & 1.6 & 5.0 \\
Hybrid & v2 & v2 & 3.40 & 67.5 & 68.5$\pm$0.2 & 82.9$\pm$0.7 & 13.7 & 19.8 \\
\bottomrule
\end{tabular}
\par\smallskip
\begin{minipage}{\linewidth}\footnotesize Signer-independent results: validation and test top-1/top-3 (\%, mean $\pm$ sd over seeds 0--2), parameter count, ONNX size and one classifier call in Chrome on the i3-1115G4 (onnxruntime-web, wasm SIMD, 1 thread, p95; seed-0 file, fp32). The deployed v2 file is the hybrid in fp16 (Section~\ref{sec:efficiency}).\end{minipage}
\end{table*}

\ifwithcomparisons% Generated by scripts/paper_tables.py -- do not edit.
\typeout{[table comparisons]}
\begin{table}[t]
\caption{Paired model comparisons}
\label{tab:comparisons}
\centering
\footnotesize
\setlength{\tabcolsep}{2.5pt}
\begin{tabular}{lrrrc}
\toprule
Comparison & $\Delta$ top-1 & 95\% CI & $p_{\mathrm{Holm}}$ & Signers \\
\midrule
GRU vs. CNN & $-$1.5 & [$-$1.9, $-$1.1] & $<$0.001 & 0/4 \\
Transf. vs. CNN & +0.1 & [$-$0.3, 0.5] & 0.717 & 2/4 \\
\midrule
+ v2 features & +0.2 & [$-$0.2, 0.6] & 0.688 & 1/4 \\
+ v2 recipe & +1.0 & [0.6, 1.4] & $<$0.001 & 4/4 \\
+ hybrid arch. & +4.9 & [4.5, 5.3] & $<$0.001 & 4/4 \\
Hybrid vs. CNN & +6.1 & [5.6, 6.5] & $<$0.001 & 4/4 \\
\bottomrule
\end{tabular}
\par\smallskip
\begin{minipage}{\linewidth}\footnotesize Paired comparisons on the test split (percentage points of top-1; each sequence's score is its correctness averaged over the three seeds). The + rows add one change to the previous model: v2 features to the v1 CNN, then the v2 recipe, then the hybrid architecture; the last row compares the hybrid with the v1 CNN. Intervals: paired bootstrap over sequences; $p$: sign-flip permutation test, Holm-adjusted over the six rows. Signers: test signers (of four) on whom the first model is better; sequences of one signer are not independent, so this column is the more conservative evidence.\end{minipage}
\end{table}
\fi

\paragraph*{Results}
Signer independence costs a lot: on a signer-dependent split (another test set, so not a
paired comparison) the same CNN reaches \V{m.dep.cnn.top1}\% test top-1, against the
\V{m.cnn.top1}\% on unseen signers in Table~\ref{tab:models}. In paired tests over the
\V{cmp.n} test sequences with Holm correction\ifwithcomparisons\ (Table~\ref{tab:comparisons})\fi\
(differences computed before rounding),
the transformer's difference from the CNN is \V{cmp.tf.sig} and the GRU differs from it by
\V{cmp.gru.diff} points. In the v2 ablation, the v2 features did not help measurably with the CNN and the v1
recipe (\V{cmp.feat.diff} points); the v2 recipe added \V{cmp.recipe.diff} points and the hybrid
\V{cmp.arch.diff}. The hybrid step also multiplies the parameter count by \V{m.paramratio} and the
training time by \V{m.timeratio}, so its gain cannot be credited to the architecture alone. These
tests treat the sequences as independent and average over seeds, so they exclude seed and signer
variation. Per signer, the hybrid's top-1 ranges from \V{sig.d.min}\% to \V{sig.d.max}\%, and it
beats the CNN for all four test signers, which is weak evidence about new ones.

The shipping rule took the highest \emph{test} top-1 with a p95 browser latency under
\V{budget.ms} ms and a file under ten megabytes; the size cap alone excludes an ensemble of the
hybrid's three seeds (\V{m.d.ens.top1}\% at \V{m.d.ens.mb}\,MB in fp16). Because the rule used the
test split, we also checked validation top-1, which selects the same model (\V{m.d.val}\% against
at most \V{m.c.val}\% for the others). The fitted temperature is below one (\V{cal.T}): the raw
model is underconfident, consistent with label smoothing. On the test split, temperature scaling
lowers the expected calibration error (\V{cal.bins} equal-width bins) from \V{cal.test.ece0}
\V{cal.test.ece0ci} to \V{cal.test.ece1} \V{cal.test.ece1ci}. The validation threshold
(\V{cal.thr}\%) accepts \V{cal.test.cov}\% of test sequences at \V{cal.test.sel}\% selective
accuracy, close to its validation operating point (\V{cal.val.cov}\% at \V{cal.val.sel}\%).

\paragraph*{Training on a second dataset (v3)}
v3 trains the hybrid with the same recipe on the Google ISLR training split plus ASL Citizen's
training videos of the mapped signs (\V{v3.train} sequences by \V{v3.trainsigners} signers), with
augmentation for webcam conditions: random aspect ratio, frame rate and dropped hand frames. Epoch
and seed are selected by the unweighted mean of the two validation top-1 scores, and each test set
was evaluated once, after v3 was selected\ifwithvthree\ (Table~\ref{tab:v3})\fi. On ASL Citizen,
v2 reaches \V{cit.v2.top1}\% (95\% bootstrap interval over signers \V{cit.v2.top1cis}),
\V{cit.v2.drop} points below its own test split, less on signs mapped through a variant
(\V{cit.v2.variant}\%) than on exact matches (\V{cit.v2.exact}\%), and its calibration degrades: the error after scaling is \V{cit.v2.ece}, and
the threshold accepts only \V{cit.v2.cov}\% of videos. v3 reaches \V{cit.v3.top1}\%, but it has
seen ASL Citizen's training signers, so that score is in-domain and not a fair comparison with v2.
On the Google ISLR test split, v3 reaches \V{v3c.ks.b}\% against \V{v3c.ks.a}\% for the deployed
v2 file, \V{v3c.ks.diff} points\ifwithvthree\else\ \V{v3c.ks.ci} (Holm $p$ \V{v3c.ks.p})\fi,
positive for each test signer. This compares two files, a
validation-selected v3 seed and v2's seed 0, not two training procedures.

\ifwithvthree% Generated by scripts/paper_tables.py -- do not edit.
\typeout{[table v3]}
\begin{table}[t]
\caption{v2 and v3 on the three test sets}
\label{tab:v3}
\centering
\footnotesize
\setlength{\tabcolsep}{2.5pt}
\begin{tabular}{lrrrr}
\toprule
Test set & v2 & v3 & v3$-$v2 [95\% CI] & $p_{\mathrm{Holm}}$ \\
\midrule
Google ISLR & 68.4 & 69.4 & +1.0 [0.5, 1.5] & $<$0.001 \\
ASL Citizen & 43.1 & 92.1 & +49.0 [47.3, 50.7] & $<$0.001 \\
Live, browser & 34.0 & 44.7 & +10.7 [$-$1.3, 23.3] & 0.238 \\
Live, Python & 40.0 & 46.0 & +6.0 [$-$6.0, 18.0] & 0.508 \\
\bottomrule
\end{tabular}
\par\smallskip
\begin{minipage}{\linewidth}\footnotesize v2 (Google ISLR only) against v3 (Google ISLR plus ASL Citizen train, with webcam augmentation): top-1 (\%) on each test set ($n$ = \V{v3c.ks.n} sequences, \V{v3c.cit.n} videos and \V{v3c.web.n} clips), each evaluated once. Differences are paired over the same items, with a paired bootstrap interval and a sign-flip test that treat items as independent (per-signer differences are in the text); browser scores are means over three runs, and the four top-1 comparisons form one Holm family, which keeps the ASL Citizen comparison although its direction was expected. v3 has seen ASL Citizen's training signers, so its ASL Citizen score is in-domain, v2's cross-dataset.\end{minipage}
\end{table}
\fi
